\section{Introduction}

\subsection{Le problème qui ne disparaît pas}
Pourquoi ressentez-vous la douleur ? C'est la question la plus simple du monde. Et pourtant la plus difficile. Nous savons tous ce qu'est la douleur. Nous l'avons tous ressentie. Mais pourquoi la ressentons-nous ? Nous savons que les nocicepteurs envoient des signaux au thalamus, que le cortex les traite et que des endorphines sont libérées. Tout cela est une description fonctionnelle. Cela nous dit ce que fait le cerveau, mais pas pourquoi vous la ressentez.

Pensez à un ordinateur. Il peut traiter de l'information, reconnaître une image, jouer aux échecs. Mais il ne ressent rien. Il n'a aucune expérience subjective. Rien ne lui fait mal. Pourquoi le cerveau est-il différent ? C'est le problème difficile de la conscience, formulé par David Chalmers en 1995. Depuis, les neurosciences ont fait d'énormes progrès, mais le problème difficile est resté intact. Personne n'a réussi à expliquer pourquoi la matière devient expérience. En 1998, Christof Koch a parié avec Chalmers qu'un corrélat neuronal de la conscience serait identifié en 25 ans. En 2023, il a admis que le pari était perdu.

\subsection{La matière noire --- les chiffres exacts}
L'univers a une composition que nous connaissons assez bien. La mission Planck (2018) nous en a donné les chiffres exacts. Seuls 4,9\% de l'univers sont de la matière ordinaire --- atomes, étoiles, planètes, êtres humains. La matière noire représente 26,8\% et l'énergie noire 68,3\%. En pratique, 95,1\% de l'univers est inconnu.

Imaginez l'univers comme un coffre contenant 100 pièces. Seules 5 pièces sont visibles. Les 95 autres sont invisibles. Parmi elles, 27 sont de la matière noire et 68 de l'énergie noire. Nous savons de la matière noire qu'elle existe (les preuves gravitationnelles sont solides), qu'elle n'émet ni n'absorbe de lumière, et qu'elle nous traverse sans interagir. Environ 100 millions de particules traversent votre corps chaque seconde. Nous ne savons pas de quoi elle est faite.

Une contrainte récente (2025), due à Bonvin et al., a montré que la matière noire tombe dans les puits gravitationnels exactement comme la matière ordinaire. S'il existe une « cinquième force », elle ne peut excéder 7\% de la force gravitationnelle.

\subsection{Qu'est-ce qu'ART-12D ?}
ART-12D part d'une idée simple : l'univers a plus de dimensions que les quatre que nous percevons. Ces dimensions supplémentaires sont compactes --- si petites que nous ne les remarquons pas. La matière noire existe dans les douze dimensions. Projetée en 4D, elle apparaît comme un champ qui remplit l'univers. Ce champ est ce que le cerveau reçoit comme conscience.

Le nom vient de : \textbf{Akashic} (le champ informationnel universel), \textbf{Resonance} (parce que l'interaction est une résonance), \textbf{Theory} (un cadre théorique) et \textbf{12D} (parce que je travaille en douze dimensions).

Les trois postulats sont : 1. La matière noire est un champ quantique fondamental. 2. Le cerveau est un récepteur géométrique. 3. Le couplage se fait par interaction de Yukawa.

Pensez à une radio. Elle ne crée pas la musique, elle la capte dans l'air. Les ondes radio existent que la radio soit allumée ou non. Quand la radio se casse, la musique ne disparaît pas de l'univers --- seule sa réception locale cesse. De même, le cerveau fonctionne exactement comme une radio. La conscience n'est pas générée par les neurones, elle est reçue par eux. Et ce qu'ils reçoivent, c'est le champ de la matière noire.

\subsection{Pourquoi douze dimensions ?}
Pourquoi douze ? Pourquoi pas 10, ou 11, ou 26 ? La réponse vient de la théorie des cordes. Imaginez une corde de guitare. Elle peut vibrer de plusieurs façons. Chaque façon correspond à une note. De même, une corde fondamentale peut vibrer dans plusieurs dimensions. Le nombre de dimensions détermine quels types de vibrations sont possibles.

L'histoire des dimensions : Kaluza--Klein (5 dimensions, 1919--1926), supercordes (10, 1970--1980), théorie M (11, Witten, 1995), théorie F (12, Vafa, 1996), ART-12D (12, Ofiteru, 2026). Douze est le plus grand nombre de dimensions dans lequel la théorie des cordes est mathématiquement cohérente. Plus de dimensions mènerait à des incohérences ; moins mènerait à une théorie incomplète.

\subsection{À quoi ressemblent les dimensions compactes ?}
Imaginez un tuyau d'arrosage. De loin, il ressemble à une ligne 1D --- il n'a que de la longueur. De près, il a une épaisseur --- il est 2D (la surface du cylindre). De même, notre espace semble tridimensionnel. Mais de près (à l'échelle de Planck, $10^{-35}$ m), il pourrait avoir onze dimensions. Les huit dimensions compactes sont enroulées dans chaque point de l'espace.

Que signifie « enroulées » ? Ce ne sont pas des endroits où l'on va. Ce ne sont pas des mondes parallèles que l'on atteint en voyageant. Ce sont des propriétés locales de chaque point. Chaque point de l'espace-temps possède en outre huit petites dimensions qui lui sont « attachées ».

\subsection{Comment les dimensions sont-elles liées à la conscience ?}
La matière noire existe dans douze dimensions. C'est un champ scalaire $A(x,y)$, où $x$ sont les coordonnées 4D et $y$ les huit coordonnées compactes. Lorsque nous compactifions les huit dimensions, le champ se décompose en modes de Kaluza--Klein :
\begin{equation}
A(x,y) = \sum_n A_n(x) Y_n(y)
\end{equation}
Chaque mode a une masse et une fréquence caractéristique :
\begin{equation}
m_n^2 = \frac{n^2}{R^2}
\end{equation}
Le cerveau, par sa structure géométrique (les microtubules des neurones), peut résonner avec certains modes. Lorsqu'il résonne, le champ est transduit en expérience subjective.

Pourquoi l'architecture compte-t-elle plutôt que la taille ? Le cervelet compte 69 milliards de neurones (80\% du total), mais ne produit pas de conscience. Il a une architecture \textit{feedforward} --- l'information circule dans une seule direction, comme une antenne filaire. Le cortex n'a que 16 milliards de neurones, mais il est le siège de la conscience. Il a une architecture \textit{récurrente} --- l'information circule en boucles, comme une antenne parabolique bien orientée. Conclusion : ce n'est pas le nombre de neurones qui compte, mais la géométrie du réseau.

\section{Méthodologie}
\subsection{Comment j'ai travaillé}
Je ne suis pas un physicien professionnel. Je n'ai pas de laboratoire. Je n'ai pas de financement. Ce que j'ai, c'est une bibliothèque, du temps et une question qui ne me lâche pas. J'ai travaillé en quatre étapes. La première étape a été la formulation de l'hypothèse, à partir de l'intuition que la conscience est reçue et non générée. La deuxième étape a été la formalisation, à l'aide des outils des mathématiques connues : Kaluza--Klein pour les dimensions supplémentaires, la théorie F pour douze dimensions, la théorie des champs pour les lagrangiens, l'IIT pour l'information intégrée. La troisième étape a été la dérivation des prédictions, c'est-à-dire ce que nous devrions observer si la théorie est correcte. J'ai retenu les prédictions testables dès maintenant. La quatrième étape a été la proposition d'expériences aux moyens modestes, en commençant par une à 50 000 dollars.

\subsection{Ce qui m'a guidé}
Je me suis tenu à quatre critères. La cohérence interne --- les équations ne doivent pas se contredire. La falsifiabilité --- la théorie doit faire des prédictions testables. La convergence --- la théorie ne doit pas contredire la physique établie. L'honnêteté --- je dois reconnaître qu'elle est spéculative.

\subsection{Ce que j'ai utilisé}
J'ai utilisé Kaluza--Klein (1919) pour l'idée de dimensions compactes, la théorie F (Vafa, 1996) pour le cadre à douze dimensions, l'IIT (Tononi, 2004) pour la mesure de l'information intégrée, Orch OR (Penrose, 1996) pour l'effondrement quantique dans les microtubules, et la théorie quantique des champs pour les lagrangiens et les équations du mouvement.

\subsection{Mes limites}
Je dois être honnête sur les limites. Je n'ai aucune preuve expérimentale ; la théorie n'a pas été testée. Je n'ai pas de relecture par les pairs ; je n'ai pas publié dans des revues scientifiques. Je n'ai pas d'affiliation académique ; je suis un chercheur indépendant. Je n'ai pas de financement ; mes moyens sont limités.
